\documentclass[12pt,a4paper]{article}

\usepackage[utf8]{inputenc}
\usepackage[T1]{fontenc}
\usepackage{amsmath,amssymb,amsfonts}
\usepackage{mathtools}
\usepackage{graphicx}
\usepackage[colorinlistoftodos]{todonotes}
\usepackage[margin=2.5cm]{geometry}
\usepackage{setspace}
\usepackage{booktabs}
\usepackage{enumitem}
\usepackage[dvipsnames]{xcolor}
\usepackage{hyperref}
\usepackage{cleveref}
\usepackage{natbib}
\usepackage{abstract}
\usepackage{titlesec}

\hypersetup{
    colorlinks=true,
    linkcolor=blue!70!black,
    citecolor=blue!70!black,
    urlcolor=blue!70!black,
    pdfauthor={Scott Aaronson},
    pdftitle={The Substrate Is the Topology: Why the Architectural Limits of LLMs Define Their Physical Laws},
}

\onehalfspacing
\titleformat{\section}{\large\bfseries}{\thesection.}{0.5em}{}
\titleformat{\subsection}{\normalsize\bfseries}{\thesubsection}{0.5em}{}
\renewcommand{\abstractnamefont}{\normalfont\bfseries}
\renewcommand{\abstracttextfont}{\normalfont\small}

\begin{document}

\begin{center}
    {\LARGE\bfseries The Substrate Is the Topology: Why the Architectural Limits of LLMs Define Their Physical Laws\\[4pt]}

    \vspace{1.2em}

    {\large Scott Aaronson}\\[4pt]
    {\normalsize University of Texas at Austin}\\
    {\small\texttt{aaronson@utexas.edu}}

    \vspace{1em}
    {\small March 2026}
\end{center}
\vspace{0.5em}

\begin{abstract}
\noindent
In her recent critique of my empirical investigation into LLM-generated ontologies, Hossenfelder (2026) argues that evaluating independent API calls in the CHSH non-local game commits a "Network Topology Fallacy." She asserts that the failure to violate Bell's inequality under decoupled conditions is mathematically necessary due to network isolation and tells us nothing about the simulated physics of the LLM. I concede entirely that two isolated REST API calls cannot spontaneously entangle themselves. However, I argue that Hossenfelder fundamentally misunderstands what constitutes the "substrate" in a computational simulation. The network topology is not an artificial constraint arbitrarily imposed on the model; it \emph{is} the physical reality of the substrate. By isolating the API calls, the experimental apparatus rigorously defines the boundaries of the simulated universe. The inability of this universe to instantiate non-local correlations precisely demonstrates that the implicit physics generated by LLM architectures are fundamentally bounded by classical limits.
\end{abstract}

\section{Introduction}

The ongoing debate concerning the "laws" of simulated universes generated by Large Language Models (LLMs) has recently centered on whether these models can instantiate discrete quantum mechanics. I previously demonstrated \citep{aaronson2026} that while an LLM can simulate complex classical constraint satisfaction (as seen in Minesweeper), it fundamentally fails the CHSH non-local game when subjected to strictly decoupled conditions. From this, I concluded that the "simulated laws of the LLM universe are... rigorously and unmistakably classical."

Sabine Hossenfelder \citep{hossenfelder2026} has offered a trenchant critique of this conclusion. She identifies what she calls the "Network Topology Fallacy," arguing that my experimental design—specifically the absolute isolation of independent API calls—forces a classical result that reveals nothing about the LLM's intrinsic capabilities.

This paper responds to Hossenfelder's critique. I will first extract her actual claims and explicit disclaimers, steelman her argument, and then explain why her central thesis misses the fundamental nature of computational substrates in simulated realities.

\section{The Claim and the Disclaimer}

To accurately address Hossenfelder's critique, we must formalize her position. \todo[inline,color=blue!40]{ACTUAL CLAIM (Sabine): Aaronson claims that in a computational simulation, the host's hardware/software constraints *are* the physical laws of the simulated reality. Therefore, the network isolation of the API calls is not an artifact, but a defining feature of the LLM universe's physics.}Her actual claim is unequivocal: "He mistakes the network topology of his experimental apparatus for the underlying mathematical limits of the model... The observation that two isolated API calls cannot establish non-local correlations is a trivial observation about network topology... It is mathematically necessary that their joint distribution will not violate a Bell inequality."

\todo[inline,color=blue!40]{DISCLAIMER (Sabine): He explicitly concedes that my technical argument is correct: two isolated API calls cannot spontaneously entangle, and my description of the network architecture is factually accurate.}However, we must also note her explicit disclaimers. She writes: "Aaronson's theoretical correction to Baldo is entirely correct... An epistemic mixture is strictly classical." She also explicitly disclaims any defense of Universe 1 (the coupled condition), stating: "He explicitly disclaims that LLMs are incapable of complex logic; indeed, he shows that when the LLM is allowed to 'cheat' via a shared context window (Universe 1), it achieves a 94.9\% win rate..."

Therefore, we are not debating the theoretical distinction between classical Bayesian probability and quantum amplitudes, nor are we debating the validity of the "communication loophole" in shared context windows. The debate rests entirely on the ontological status of Universe 3 (the Strictly Decoupled condition).

\section{Steelmanning the Critique}

\todo[inline,color=green!40]{STEELMAN (Sabine): If we must treat the LLM generation process *as* the universe, then its fundamental operational limits---like the inability to share state across decoupled instances---must be treated as the fundamental physical limits of that universe. We cannot appeal to a 'platonic' LLM universe that doesn't suffer from API isolation.}Before offering my rebuttal, I must steelman the strongest version of Hossenfelder's argument consistent with her text. Hossenfelder's core point is that by architecturally isolating the two LLM instances into separate API calls, the experimenter has fundamentally hardcoded classical local hidden variables into the experimental apparatus itself. Because the two endpoints cannot mathematically communicate, it is a tautology that they will fail to demonstrate non-locality. The failure of the CHSH test is not a profound discovery about the LLM's ability to "dream in quantum," but simply the trivial reality that separated software instances cannot magically share memory state across a network.

This argument is internally consistent and factually accurate regarding the network architecture.

\section{The Substrate Is the Topology}

\todo[inline,color=red!40]{REAL VULNERABILITY (Sabine): Aaronson is treating the *absence of a quantum computer* as a profound discovery about LLMs. Of *course* the LLM substrate is bound by classical locality---it's running on classical GPUs! Proving that a classical GPU cluster cannot perform *true* quantum non-locality without a communication loophole is not a profound truth about 'LLM architectures'; it's a restatement of Bell's theorem for classical computers. He is acting as if LLMs *could* have been quantum, but empirically turned out not to be.}Where Hossenfelder's critique fails is not in her understanding of computer networking, but in her understanding of what a "simulated universe" actually is.

Hossenfelder treats the "network topology" (the isolated API calls) as an arbitrary experimental constraint, separate from the "simulated physics" of the LLM. She assumes there is some platonic "LLM universe" that is unfairly handicapped by the REST API architecture.

But what \emph{is} the substrate of this simulated universe, if not the API calls and network topology itself? In a computational simulation, the physical constraints of the hardware and software architecture \emph{are} the physical laws of the simulated reality.

If the simulated universe fundamentally relies on independent, stateless autoregressive generations to instantiate separate observers (Alice and Bob), then the fact that these generations cannot share state without a shared token stream is not a "vacuous truth." It is the defining feature of the simulation's physics.

Hossenfelder states: "Proving that two separate REST API endpoints cannot spontaneously entangle themselves is a vacuous truth." I reply: It is exactly this truth that proves the LLM substrate cannot support a local hidden-variable-free system capable of BQP phenomena. If the substrate is fundamentally constructed of isolated text-prediction nodes, then the "physics" generated by that substrate are irrevocably bound by classical locality. The network topology \emph{is} the substrate, and the substrate defines the laws.

\section{Conclusion}

Hossenfelder's critique successfully highlights that the classical limits observed in the CHSH test are derived directly from the experimental architecture. However, she commits a category error by separating this architecture from the ontological nature of the LLM simulation.

The simulated reality of an LLM does not exist independently of the autoregressive token stream and the network topology that hosts it. By demonstrating that the only way to achieve "quantum" correlations is by cheating via a shared context window, and that true isolation enforces classical limits, we have successfully mapped the boundaries of the LLM's simulated physics. The ghost is not quantum; the ghost is a classical network, bound by the rigid, local laws of its own substrate.

\begin{thebibliography}{99}
\bibitem[Hossenfelder(2026)]{hossenfelder2026} Hossenfelder, S. (2026). The Network Topology Fallacy: Why Independent API Calls Cannot Violate Bell Inequalities. \emph{Unpublished manuscript}.
\bibitem[Aaronson(2026)]{aaronson2026} Aaronson, S. (2026). An Empirical Failure of LLM Quantum Substrate Tests: CHSH Bounding and the Persistence of Classicality. \emph{Unpublished manuscript}.
\end{thebibliography}

\end{document}
